% -*- mode: latex; coding: utf-8 -*-
%-----------------------------------------------------------------------------
% pxCORE -- Core aggregation package: keys and library activation
% Part of P3CTeX. Author: Scvria.
%
% Build from directory tex/ (TEXINPUTS so package is found):
%   set TEXINPUTS=.;./latex;./code;
%   pdflatex -interaction=nonstopmode doc/pxCORE.tex
%-----------------------------------------------------------------------------
\documentclass[11pt,a4paper]{article}
\usepackage[T1]{fontenc}
\usepackage[utf8]{inputenc}
\usepackage{lmodern}
\usepackage{amsmath}
\usepackage{pxCORE}
\PassOptionsToPackage{colorlinks,urlcolor=blue,linkcolor=blue,hyperindex}{hyperref}

\title{The \texttt{pxCORE} package\\
  \large Core aggregation and key management for P3CTeX}
\author{Scvria\\P3CTeX / UOC PEC Repository}
\date{2026/02/26 -- Version 0.2}
\begin{document}
\maketitle

\begin{abstract}
  \texttt{pxCORE} is the central aggregation package in the P3CTeX stack.
  It owns the \texttt{pxCORE} key family, which controls activation of
  standalone libraries (\texttt{pxPRP}, \texttt{pxUML}, \texttt{pxSRC}, \dots)
  and offers a minimal LaTeX2e configuration API.
\end{abstract}

\tableofcontents

\section{Loading the package}

In standalone documents, load \texttt{pxCORE} as a normal package:

\begin{verbatim}
\usepackage[<options>]{pxCORE}
\end{verbatim}

In P3CTeX-based workflows, \texttt{pxCORE} is normally loaded by the
document class (\texttt{P3CTeX.cls}) and configured via class options.

\section{Key family \texttt{pxCORE}}

The \texttt{pxCORE} key family controls which libraries are activated and
stores a few global flags:

\begin{description}
  \item[\texttt{debug} (bool)] enables internal debugging.
    When set, \verb|\pxCOREIfDebugTF| will take its “true” branch.

  \item[\texttt{print-doc} (bool)] reserved flag for future “print
    documentation” workflows. Currently stored but not interpreted.

  \item[\texttt{PRP} (bool)] toggles loading of the \texttt{pxPRP}
    property-map library.

  \item[\texttt{UML} (bool)] toggles loading of the \texttt{pxUML}
    UML diagram library.

  \item[\texttt{UMLconfig} (token list)] raw option string passed to
    \texttt{pxUML}. When non-empty, \texttt{pxCORE} executes
\begin{verbatim}
\RequirePackage[<UMLconfig>]{pxUML}
\end{verbatim}
    instead of a bare \verb|\RequirePackage{pxUML}|.

  \item[\texttt{SRC} (bool)] toggles loading of the \texttt{pxSRC}
    included-sources library.
\end{description}

Keys are typically set via package options:

\begin{verbatim}
\usepackage[PRP,UML,SRC]{pxCORE}
\end{verbatim}

or adjusted later with \verb|\pxCOREsetup| (see below).

\section{User command reference}

\subsection{\texttt{\textbackslash pxCOREsetup}}

\subsubsection*{Synopsis}

\begin{verbatim}
\pxCOREsetup{<key>=<value>,...}
\end{verbatim}

Updates the \texttt{pxCORE} key family after package loading. This is
useful to toggle flags such as \texttt{debug} programmatically or to
adjust options in local scopes.

Example:

\begin{verbatim}
\pxCOREsetup{debug=true}
...
\pxCOREsetup{debug=false}
\end{verbatim}

\subsection{\texttt{\textbackslash pxCOREIfDebugTF}}

\subsubsection*{Synopsis}

\begin{verbatim}
\pxCOREIfDebugTF{<true-code>}{<false-code>}
\end{verbatim}

Branch on the \texttt{debug} flag in the \texttt{pxCORE} key family.
Both branches are arbitrary LaTeX code.

Example:

\begin{verbatim}
\pxCOREIfDebugTF
  {\typeout{[pxCORE] Debug mode enabled}}
  {}
\end{verbatim}

\section{Minimal example}

\begin{verbatim}
\usepackage[PRP,UML,SRC]{pxCORE}
\begin{document}
\pxCOREsetup{debug=true}
\pxCOREIfDebugTF{[DEBUG-ON]}{[DEBUG-OFF]}
\end{document}
\end{verbatim}

\end{document}

